\Needspace{12\baselineskip}
\originalchapter{非线性信号的观测}
\originalsection{频谱、锁频与宽带}
时间序列的频谱把周期结构和尺度分离。对有限样本 $x_n$，先去均值，DFT分辨率为 $1/(N\Delta t)$，最高可辨频率为 $1/(2\Delta t)$。窗函数减少泄漏，却改变有效分辨率与功率归一化；不同Welch块的重叠和平均影响误差。

周期振荡常给基频和谐波，准周期运动出现多个不成有理比的基频及组合峰，混沌和噪声都可能产生宽带。仅有宽带谱不能区分确定性混沌与随机噪声。稳定谱峰也可能来自外部周期驱动或仪器干扰，不能直接推断内部极限环。

\originalsection{截面和延迟重构}
Poincaré截面把连续流约为回返映射。取 $x=0$ 截面必须还指定穿越方向，否则前后穿越混合两种回返。数值上按符号变化线性或高阶插值定位交点，比直接选取 $|x|<\epsilon$ 的一片厚层更接近真正截面。截面没有保存每条轨迹全部时间信息。

只有一个可测量量 $x(t)$ 时，可用 $(x(t),x(t+\tau),x(t+2\tau),\ldots)$ 重构几何。延迟太短时所有坐标几乎相同，太长时失去关联；样本量、噪声、采样率和系统维数限制结论。通用重构定理需要光滑确定性系统、一般观测函数和足够维度，本卷实际数据图不被当作该定理条件已经自动满足的证明。

\originalsection{公开对流数据}
12.006J PS7公开 \texttt{exp1.csv}、\texttt{exp2.csv}，题面署名Monique Dubois和Pierre Bergé，称它们是同一次对流实验两个不同位置的热梯度相关信号。本卷按原始字节保存数据并注明来源。第一列是样本编号，所核文件未说明采样间隔和物理标定，因此频率以每样本周期数表示，幅度保留原任意单位。

代码对两组信号分别去均值、计算一侧Welch谱并保存延迟重构；输出见题解与QA。两地信号形状和谱结构可比较，但不能把没有标定的量换算为绝对温差、Hz或具体Rayleigh数。谱和重构联合支持对流状态的有界判断，不能仅凭图宣布系统严格确定性混沌。

\originalsection{Fourier变换、有限记录与Wiener--Khinchin}
连续变换采用 $\hat x(\omega)=\int x(t)\e^{-i\omega t}\dd t$、逆变换含 $1/(2\pi)$。有限离散记录用 $X_k=\sum_{n=0}^{N-1}x_n\e^{-2\pi ikn/N}$，逆变换含 $1/N$。Parseval给 $\sum|x_n|^2=N^{-1}\sum|X_k|^2$，是一侧谱归一化的重要核对。

平稳过程自相关 $C(\tau)=\avg{x(t)x(t+\tau)}$ 的Fourier变换为功率谱。有限时段的周期图很噪，Welch分块窗口平均降低方差，但不增加单块的本征频率分辨。把零填充当作多出真实分辨率，是对插值与观测信息的混淆。

矩形记录窗的频谱为sinc，信号频率未落在DFT格点时泄漏到邻点；窗函数削旁瓣但扩大主瓣。非正弦周期信号的谐波属于真实波形，不应误判为多独立振子。准周期两个非公度频率给整数组合 $m\omega_1+n\omega_2$；有限观测可能无法分辨很接近的有理比。

\originalsection{Poincar\'e截面与回返时间}
自治流穿过横截面，首次正向回返定义映射 $P$。连续流的一个稳定周期轨道在截面是稳定固定点；二维环面准周期流在截面为一闭曲线；奇异吸引子可给分形点集。截面应横切流而不切向，否则回返时间与数值定位病态。

映射删掉每次回返间的时间，若研究频率/间歇必须另保存 $\tau_n$。周期强迫的每驱动周期采样是另一个“频闪截面”，相位固定而无需检测几何过零。两种取样都降维，但条件不同。

\originalsection{延迟重构的选择与局限}
从单一观测构建 $\mathbf y_n=(x_n,x_{n+\tau},\ldots,x_{n+(m-1)\tau})$。太小 $\tau$ 使坐标几乎共线，太大则混合无关时段；可用相关降到一定尺度或互信息首小值作为候选，再比较结果。嵌入维数需足够高以避免投影交叉，假近邻方法也要考虑噪声和时间相邻样本排除。

典型光滑系统的嵌入定理给一般观测在足够维度中的拓扑信息，但不保证少样本噪声数据可以精确求吸引子维数。一个三维延迟云看似厚，可能来自随机噪声或未观测驱动；本文两组公开数据只按实际支持的频谱/重构结果作判断。

\originalsection{对流数据的可复核观察}
两文件各15000点，所选一侧Welch谱主峰约 $.0629883$ 与 $.0317383$ 周期/样本。频率比接近二，但来自同次实验不同位置，可能是空间模式选择或谐波差异，不能当作控制参数逐渐改变产生倍周期的直接证据。绘图使用原始幅度与样本横轴，并把去均值、块长和重叠规则保存。

比较原始时间序列、自相关、谱和延迟图能提出模型候选，但稳定周期、准周期、噪声驱动与确定性混沌的区别需要额外证据。原文件没有实际时间步和温度校准，不写Hz、开尔文或Rayleigh数；这是真实观测限制，而非通过猜测可以补全的空白。
